\documentclass[10pt,a4paper]{amsart}
\usepackage{amsmath,amssymb,amsthm,graphicx}
% Retained preamble for the Stacks Project chapter divisors.tex (Divisors).
% The chapter file pinned-original-source.tex carries only \input{preamble} on its first line, so no
% preamble text lives inside the chapter itself.  The Stacks Project uses one single shared preamble for
% every chapter; the snapshot stages it verbatim in sources/stacks-chow/preamble.tex.  The block below is
% copied mechanically, line for line, from that file's lines 14-18 (the Xy-pic setup) and lines 154-217
% (the native theorem declarations together with the mathematical macros).  Nothing was added, removed or
% redefined.  The wrapper class and the amsmath/amssymb/amsthm/graphicx base packages are supplied by the
% assembler, so the class/\externaldocument part of the shared preamble is not repeated here.
% For commutative diagrams we use Xy-pic
\usepackage[all]{xy}

% We use 2cell for 2-commutative diagrams.
\xyoption{2cell}
% Theorem environments.
%
\theoremstyle{plain}
\newtheorem{theorem}[subsection]{Theorem}
\newtheorem{proposition}[subsection]{Proposition}
\newtheorem{lemma}[subsection]{Lemma}

\theoremstyle{definition}
\newtheorem{definition}[subsection]{Definition}
\newtheorem{example}[subsection]{Example}
\newtheorem{exercise}[subsection]{Exercise}
\newtheorem{situation}[subsection]{Situation}

\theoremstyle{remark}
\newtheorem{remark}[subsection]{Remark}
\newtheorem{remarks}[subsection]{Remarks}

\numberwithin{equation}{subsection}

% Macros
%
\def\lim{\mathop{\mathrm{lim}}\nolimits}
\def\colim{\mathop{\mathrm{colim}}\nolimits}
\def\Spec{\mathop{\mathrm{Spec}}}
\def\Hom{\mathop{\mathrm{Hom}}\nolimits}
\def\Ext{\mathop{\mathrm{Ext}}\nolimits}
\def\SheafHom{\mathop{\mathcal{H}\!\mathit{om}}\nolimits}
\def\SheafExt{\mathop{\mathcal{E}\!\mathit{xt}}\nolimits}
\def\Sch{\mathit{Sch}}
\def\Mor{\mathop{\mathrm{Mor}}\nolimits}
\def\Ob{\mathop{\mathrm{Ob}}\nolimits}
\def\Sh{\mathop{\mathit{Sh}}\nolimits}
\def\NL{\mathop{N\!L}\nolimits}
\def\CH{\mathop{\mathrm{CH}}\nolimits}
\def\proetale{{pro\text{-}\acute{e}tale}}
\def\etale{{\acute{e}tale}}
\def\QCoh{\mathit{QCoh}}
\def\Ker{\mathop{\mathrm{Ker}}}
\def\Im{\mathop{\mathrm{Im}}}
\def\Coker{\mathop{\mathrm{Coker}}}
\def\Coim{\mathop{\mathrm{Coim}}}

% Boxtimes
%
\DeclareMathSymbol{\boxtimes}{\mathbin}{AMSa}{"02}

%
% Macros for moduli stacks/spaces
%
\def\QCohstack{\mathcal{QC}\!\mathit{oh}}
\def\Cohstack{\mathcal{C}\!\mathit{oh}}
\def\Spacesstack{\mathcal{S}\!\mathit{paces}}
\def\Quotfunctor{\mathrm{Quot}}
\def\Hilbfunctor{\mathrm{Hilb}}
\def\Curvesstack{\mathcal{C}\!\mathit{urves}}
\def\Polarizedstack{\mathcal{P}\!\mathit{olarized}}
\def\Complexesstack{\mathcal{C}\!\mathit{omplexes}}
% \Pic is the operator that assigns to X its picard group, usage \Pic(X)
% \Picardstack_{X/B} denotes the Picard stack of X over B
% \Picardfunctor_{X/B} denotes the Picard functor of X over B
\def\Pic{\mathop{\mathrm{Pic}}\nolimits}
\def\Picardstack{\mathcal{P}\!\mathit{ic}}
\def\Picardfunctor{\mathrm{Pic}}
\def\Deformationcategory{\mathcal{D}\!\mathit{ef}}
\usepackage{cleveref}
\crefname{originallemma}{original lemma}{original lemmas}
\Crefname{originallemma}{Original lemma}{Original lemmas}
\makeatletter
\newlabel{coherent-lemma-flat-base-change-cohomology}{{02KH}{source}{}{sourcepointer-1}{}}
\newlabel{coherent-lemma-flat-base-change-cohomology@cref}{{[originallemma][0][]02KH}{[1][0][]source}{}{}{}}
\makeatother
\crefname{originallemma}{original lemma}{original lemmas}
\Crefname{originallemma}{Original lemma}{Original lemmas}
\makeatletter
\newlabel{schemes-lemma-push-forward-quasi-coherent}{{01LC}{source}{}{sourcepointer-2}{}}
\newlabel{schemes-lemma-push-forward-quasi-coherent@cref}{{[originallemma][0][]01LC}{[1][0][]source}{}{}{}}
\makeatother
\title{A point outside the image of a quasi-compact and quasi-separated morphism is not weakly associated to the pushforward}
\author{The Stacks Project Authors}
\date{stacks-project:divisors.tex:local-snapshot-20261012; GNU FDL 1.2 or later}
\begin{document}
\maketitle
\noindent\textit{Editorial note.} All delivered mathematical text is the author's own, copied from the pinned Stacks Project chapter file divisors.tex; no statement or proof step was rewritten, reordered or removed.\par\medskip
\noindent\textit{Editorial note.} The canonical Stacks tag of the counted statement is 0AVN (divisors-lemma-weakass-pushforward). Following the multi-claim convention for this source, source\_url is the source-level URL https://stacks.math.columbia.edu/ and the tag is carried in claim\_key and here rather than in a per-claim tag URL.\par\medskip
\noindent\textit{Editorial note.} Two slips of the author inside the counted proof are reproduced verbatim: the displayed injectivity uses M\_\{a\_r\} where the preceding cover is indexed by a\_1, ..., a\_n, and the element m is said to lie in the stalk M\_p where p should be the maximal ideal m of the local ring. No mathematical text was corrected.\par\medskip
\noindent\textit{Editorial note.} The definition of weakly associated point is attached as a context span because the statement uses the term without repeating the definition.\par\medskip
\noindent\textit{Editorial note.} The preamble is the chapter's shared Stacks Project preamble, retained as sources/stacks-divisors/retained-preamble.tex: a line-for-line excerpt of sources/stacks-chow/preamble.tex lines 14-18 and 154-217. No macro was invented or redefined.\par\medskip
\noindent\textit{Allowed standard prerequisites (editorial).}\begin{itemize}
\item Flat base change for quasi-coherent cohomology of a quasi-compact and quasi-separated morphism. Cohomology of Schemes, Lemma labelled coherent-lemma-flat-base-change-cohomology (canonical Stacks tag 02KH); invoked at proof line 736. Supplies the identification (g\textasciicircum{}* f\_* F)\_s = (f'\_* (g')\textasciicircum{}* F)\_s, which lets the proof assume that the base is local.
\item The pushforward of a quasi-coherent module under a quasi-compact and quasi-separated morphism is quasi-coherent. Schemes, Lemma labelled schemes-lemma-push-forward-quasi-coherent (canonical Stacks tag 01LC); invoked at proof line 741. Lets the author write f\_* F as the sheaf associated to a single A-module M, which is the object the rest of the argument estimates.
\item A prime is weakly associated to a module exactly when it is the radical of the annihilator of some element. Algebra, Definition labelled algebra-definition-weakly-associated and Algebra, Lemma labelled algebra-lemma-weakly-ass-local (canonical Stacks tags 0547 and 0566), recalled in the section introducing the definition. Turns the conclusion "s is not weakly associated to f\_* F" into the explicit statement that no nonzero element m of the stalk has all a\_i\textasciicircum{}n m equal to zero.
\end{itemize}\medskip
\ifdefined\hypertarget\hypertarget{sourcepointer-1}{}\fi
\noindent\textit{Source reference.} Original lemma 02KH: sources/stacks-chow/tags.txt, source lines 3372--3372. Cross-chapter pointer invoked at proof line 736 for flat base change; the resulting identification of stalks is then used as the reduction to a local base, which the counted proof carries out itself.\par\medskip
\ifdefined\hypertarget\hypertarget{sourcepointer-2}{}\fi
\noindent\textit{Source reference.} Original lemma 01LC: sources/stacks-chow/tags.txt, source lines 2081--2081. Cross-chapter pointer invoked at proof line 741 for quasi-coherence of f\_* F; the module M and every estimate on it are then produced inside the counted proof.\par\medskip
\begin{lemma}
\label{lemma-weakass-pushforward}
Let $f : X \to S$ be a quasi-compact and quasi-separated morphism of schemes.
Let $\mathcal{F}$ be a quasi-coherent $\mathcal{O}_X$-module.
Let $s \in S$ be a point which is not in the image of $f$. Then
$s$ is not weakly associated to $f_*\mathcal{F}$.
\end{lemma}

\begin{proof}
Consider the base change $f' : X' \to \Spec(\mathcal{O}_{S, s})$
of $f$ by the morphism $g : \Spec(\mathcal{O}_{S, s}) \to S$
and denote $g' : X' \to X$ the other projection.
Then
$$
(f_*\mathcal{F})_s = (g^*f_*\mathcal{F})_s = (f'_*(g')^*\mathcal{F})_s
$$
The first equality because $g$ induces an isomorphism on local
rings at $s$ and the second by flat base change (Cohomology of Schemes, Lemma
\ref{coherent-lemma-flat-base-change-cohomology}). Of course
$s \in \Spec(\mathcal{O}_{S, s})$ is not in the image of $f'$.
Thus we may assume $S$ is the spectrum of a local ring
$(A, \mathfrak m)$ and $s$ corresponds to $\mathfrak m$.
By Schemes, Lemma \ref{schemes-lemma-push-forward-quasi-coherent}
the sheaf $f_*\mathcal{F}$ is quasi-coherent, say corresponding
to the $A$-module $M$. As $s$ is not in the image of $f$ we see that
$X = \bigcup_{a \in \mathfrak m} f^{-1}D(a)$ is an open covering.
Since $X$ is quasi-compact we can find $a_1, \ldots, a_n \in \mathfrak m$
such that $X = f^{-1}D(a_1) \cup \ldots \cup f^{-1}D(a_n)$. It follows
that
$$
M \to M_{a_1} \oplus \ldots \oplus M_{a_r}
$$
is injective. Hence for any nonzero element $m$ of the stalk $M_\mathfrak p$
there exists an $i$ such that $a_i^n m$ is nonzero for all $n \geq 0$.
Thus $\mathfrak m$ is not weakly associated to $M$.
\end{proof}

\begin{definition}
\label{definition-weakly-associated}
Let $X$ be a scheme.
Let $\mathcal{F}$ be a quasi-coherent sheaf on $X$.
\begin{enumerate}
\item We say $x \in X$ is {\it weakly associated} to $\mathcal{F}$
if the maximal ideal $\mathfrak m_x$ is weakly associated to the
$\mathcal{O}_{X, x}$-module $\mathcal{F}_x$.
\item We denote $\text{WeakAss}(\mathcal{F})$ the set of weakly associated
points of $\mathcal{F}$.
\item The {\it weakly associated points of $X$} are the weakly associated
points of $\mathcal{O}_X$.
\end{enumerate}
\end{definition}

\end{document}
